\documentclass{article}
\usepackage[UTF8]{ctex}
\usepackage{amsmath,amssymb,amsfonts}
\usepackage{enumitem}
\usepackage{geometry}
\geometry{a4paper, margin=1in}
% === 错题本排版优化（不参与题目标识） ===
\setlength{\parindent}{0pt}
\setlength{\parskip}{0.35em plus 0.10em minus 0.08em}
\setlength{\emergencystretch}{3em}
\setlist{itemsep=0.25em, topsep=0.40em, parsep=0pt}
\allowdisplaybreaks[2]
% === 错题本排版优化结束 ===
\begin{document}
\section*{长安大学2026年数学分析真题}
\section*{186．长安大学2026年数学分析试卷}

\subsection{题目1}
1．计算极限 \(\lim _{x \rightarrow 0} \frac{\cos (\sqrt{2} x)-e^{-x^{2}}+\frac{x^{4}}{3}}{x^{6}}\) ．

\subsection{题目2}
2．设函数 \(y=f(x)\) 由方程 \(y^{3}+x y^{2}+x^{2} y+6=0\) 确定，求 \(f(x)\) 的极值．

\subsection{题目3}
3．计算积分 \(J=\int_{0}^{1} \ln \frac{1+x}{1+x^{2}} \mathrm{~d} x\) ．

\subsection{题目4}
4．计算重积分 \(\iint_{D} y \sin (x+y) \mathrm{d} x \mathrm{~d} y\) ，其中 \(D=\left\{(x, y) \left\lvert\, 0 \leq x+y \leq \frac{\pi}{2}\right., 0 \leq x-y \leq \frac{\pi}{2}\right\}\) 。

\subsection{题目5}
5．已知有向曲线 \(L\) 为球面 \(x^{2}+y^{2}+z^{2}=2 x\) 与平面 \(2 x-z-1=0\) 的交线，从 \(z\) 轴正向往 \(z\) 轴负向看为逆时针方向，计算曲线积分

\[
\oint_{L}\left(6 x y z-y z^{2}\right) \mathrm{d} x+2 x^{2} z \mathrm{~d} y+x y z \mathrm{~d} z .
\]

\subsection{题目6}
6．设 \(S\) 为有界数集，证明：若 \(\sup S=a \notin S\) ，则存在严格单调递增的数列 \(\left\{x_{n}\right\} \subset S\) ，使得

\[
\lim _{n \rightarrow \infty} x_{n}=a
\]

\subsection{题目7}
7．设函数 \(f(x), g(x)\) 在 \([a, b]\) 上连续，在 \((a, b)\) 内具有二阶导数且存在相等的最大值，\(f(a)=g(a)\) ， \(f(b)=g(b)\) ，证明：存在 \(\xi \in(a, b)\) ，使得 \(f^{\prime \prime}(\xi)=g^{\prime \prime}(\xi)\) 。

\subsection{题目8}
8．证明：若 \(f(x)\) 在 \([a, b]\) 上可积，\(F(x)\) 在 \([a, b]\) 上连续，且除有限个点外有 \(F^{\prime}(x)=f(x)\) ，则有

\[
\int_{a}^{b} f(x) \mathrm{d} x=F(b)-F(a)
\]

\subsection{题目9}
9．证明：若 \(f(x)\) 在 \([a,+\infty)\) 上可导，且 \(\int_{a}^{+\infty} f(x) \mathrm{d} x\) 与 \(\int_{a}^{+\infty} f^{\prime}(x) \mathrm{d} x\) 都收敛，则 \(\lim _{x \rightarrow+\infty} f(x)=0\) ．

\subsection{题目10}
10．解答如下问题：设 \(f(x)=\sum_{n=2}^{\infty} \frac{x^{n}}{n^{2} \ln n}\) 。证明：

\begin{itemize}
  \item[（1）] 证明：\(f(x)\) 在 \([-1,1]\) 上连续．
  \item[（2）] 证明： \(\lim _{x \rightarrow 1^{-}} f^{\prime}(x)=+\infty\) ．
  \item[（3）] 证明： \(\lim _{x \rightarrow 1^{-}} f^{\prime}(x)=+\infty\) ．
  \item[（4）] 证明：\(f(x)\) 在 \(x=1\) 处左不可导．
\end{itemize}

\subsection{题目11}
11．解答如下问题：设 \(I(k, a)=\int_{0}^{+\infty} e^{-k x} \frac{\sin (a \dot{x})}{x} \mathrm{~d} x\) ．\\
（1）证明：对固定的 \(k \in[0,+\infty)\) ，积分 \(I(k, a)\) 关于 \(a\) 在 \(|a| \geq \delta(\delta>0)\) 上一致收敛，对固定的 \(a \neq 0\) ，积分 \(I(k, a)\) 关于 \(k\) 在 \([0,+\infty)\) 上一致收敛。\\
（2）计算积分 \(I=\int_{0}^{+\infty} \frac{\sin x}{x} \mathrm{~d} x\) 的值．
\end{document}
